\documentclass[11pt]{article}


\usepackage[inline]{asymptote}
\usepackage{asypictureB}
%\usepackage[default]{frcursive}
\usepackage[a4paper,text={17.5cm,25.4cm},centering,top=2.1cm]{geometry} 
\usepackage{amsfonts}
%\usepackage{amsmath}
\usepackage{amssymb,amsmath,stackrel,mathrsfs} %wasysym
\usepackage{multicol}
\usepackage{pdflscape}% girar una pagina, hay que usar \begin{landscape}En este punto inicia la hoja en vertical \end{landscape}
%\usepackage[mathcal]{euscript}
%\usepackage{amssymb}
\usepackage{esvect}%para la flecha de vector 

\usepackage[T1]{fontenc}
\usepackage[spanish,es-noshorthands]{babel}
%\usepackage[latin1,ansinew]{inputenc}\usepackage{calligra}
%\usepackage[pdftex]{graphicx}
%\usepackage{etex}
\usepackage{calligra}
\usepackage{array}
 \usepackage{booktabs}
 \usepackage{ifsym}
\usepackage{pgf,tikz,pgfplots}
\usepackage{tkz-graph}%paquete (di)grafos
\usepackage{tkz-euclide}% Para realizar diagramas que coloreen conjuntos
\usetikzlibrary[decorations.text]% Decorar curvas
%\usetikzlibrary{matrix}
\usetikzlibrary{arrows,matrix}
%\usetkzobj{all}

\pgfplotsset{compat=1.15}
\usepackage{mathrsfs}
\usetikzlibrary{arrows}
\usepackage{makecell} %Paquete para controlar altura de celdas
\usepackage{verbatim} 
\usetikzlibrary{babel}%para no general error al usar babel y tikz simultaneamente.
\usetikzlibrary{patterns,patterns.meta}%%%%%%%%%%%para  pintar de forma rrallada un rectangulo

\usepackage{pictex}
\usepackage{verbatim}%me sirve para usar comentarios con mucha lineas 
\usepackage{amstext}
\usepackage{enumerate}
%\usepackage{enumitem}
\usepackage{amsthm}
\usepackage{xcolor}
\usepackage{mathdots}% para tener mas opciones de puntos suspensivos

%\usepackage{pgf,tikz}%,pgfplots
%\pgfplotsset{compat=1.15}
\usepackage{mathrsfs}
%\usetikzlibrary{arrows}
%\usepackage{etex,tkz-euclide} 
%\usetkzobj{all}
\usepackage{color}

\usepackage{cancel}
\newcommand{\Cancelar}[2][black]{{\color{#1}\cancel{\color{black}#2}}}%%% para cancelar con colores recordar usar mayusculas
\usepackage{lscape}%paquete para girar hojas


\usepackage{bm}









\newcommand{\dis}{\displaystyle}
\newcommand{\R}{\mathbb{R}}
\newcommand{\re}{\mathbb{R}}
\newcommand{\I}{\mathbb{I}}
\newcommand{\Q}{\mathbb{Q}}
\newcommand{\Z}{\mathbb{Z}}
\newcommand{\C}{\mathbb{C}}
\newcommand{\N}{\mathbb{N}}
\newcommand{\K}{\mathbb{K}}
%\newcommand{\B}{\mathcal{B}}
\newcommand{\V}{\mathbb{V}}
\newcommand{\can}{\mathcal{C}}
\newcommand{\M}{\mathbb{M}}
\newcommand{\Pol}{\mathcal{P}}

\newcommand{\iv}{\bm{\check{\imath}}}
\newcommand{\jv}{\bm{\check{\jmath}}}
\newcommand{\kv}{\check{\bm k}}
\newcommand{\moduloprop}{\hookrightarrow_{\hspace{-2.5mm}\mbox{\tiny$^{_\neq}$}}}

\newtheorem{ejer}{Ejercicio}[]

\newtheorem{teo}{Teorema}%[section]
\newtheorem{lema}[teo]{Lema}
\newtheorem{propo}[teo]{Proposición}
\newtheorem{defi}[teo]{Definición}
\newtheorem{ejem}[teo]{}
\newtheorem*{ejemsinn}{}
\newtheorem{obs}[teo]{Observación}%[section]
\newtheorem{acla}{Aclaración}%[section]
\newtheorem{act}{Actividad}%%%%%%%%%%%%%%%%
\newtheorem*{defsinn}{{\bf Definición}:}
\newtheorem*{ejemplosinn}{{\bf Ejemplo}:}
\newtheorem*{proposinn}{{\bf\sc Proposición}}
\newtheorem*{obssinn}{Observación:}
\newtheorem*{corosinn}{Corolario}
\newtheorem*{propsinn}{Propiedades:}
%\newenvironment{proof} { \noindent {\bf Demostración: }\rm}{$\quad \hfill \blacksquare $}
\usepackage{setspace} % (es una alternativa a \renewcommand{\baselinestretch}{1.3}) usando \doublespacing \onehalfspace \singlespace \spacing{1.5}  se va cambiando segun uno quiera
\spacing{1.3}

%\newenvironment{proof}[Demostración:]
\newenvironment{sol}{ \noindent {\bf Solución:} \rm}{%\hfill \protect\tikz[overlay,yshift=-4pt]\protect\node[xshift=8mm, yshift=13mm, blue]() {\huge $\checkmark$ };
}

%%%%%%%%%%%%%%%%%%%%%%%%%%%%%%%%%%%%%%%%


%\renewcommand{\baselinestretch}{1.5}
%%%%%%%%%%%%%%%%%%%%%%%%%%%%%%%%%%%%%%%%%%%%%%%%%%%%%%%%%%%%%%%%%%%%%%%%%%%%%%%%%%%%%%%%%%%%%%%%%%%%%%%%%%%%%%%%%%%%%%%%%%%%%%%%%%%%%%%%%%%%%%%%%%%%%%%%%%%%%%%%%%%%%%%%%%%%%%%%%%%%%%%%%%%%%%%%%%%%%%%%%%%%%%%%%%%%%%%%%%%%%%%%%%%%%%%%%%%%%%%%%%%%%%%%%%%%%%%%%%%%%%%%%%%%%%%%%%%%%%%%
\usepackage[framemethod=TikZ]{mdframed}%Paquete para estilos de los bordes de los ambientes Importante este paquete (mdframed) define un entorno  que trata automáticamente los saltos de página en el texto enmarcado.
            \mdfsetup{skipabove=1\baselineskip, skipbelow=1.5\baselineskip ,             frametitlefont= \sffamily\bfseries , needspace=4\baselineskip , splittopskip=1.5\baselineskip} %
            \mdfsetup{apptotikzsetting={\tikzset{mdfbackground/.append style={draw=none}}}} 
%%%%%%%%%%%%%%%%%%%%%%%%%%%%%%%%%%%%%%%%%%%%%%%%%%%%%%%%%%%%%%
%
%Estilo para Cuentas auxiliares
%
%%%%%%%%%%%%%%%%%%%%%%%%%%%%%%%%%%%%%%%%%%%%%%%%%%%%%%%%%%%%%%
\mdfdefinestyle{Cuentaauxi}{%skipabove=10pt, skipbelow=0,
userdefinedwidth=1.02\textwidth,align=center,
  everyline=true,
  ignorelastdescenders=true,
  middlelinewidth=1pt,
  linecolor=black!20,
  outerlinewidth=1pt,
  %leftline=false,topline=false
  %rightline=false,bottomline=false,
  %backgroundcolor=black!20,
  innerleftmargin=2mm, innerrightmargin=2mm,
  innerbottommargin=3mm,
  innertopmargin=1mm,
  frametitleaboveskip=2mm,
  frametitlebelowskip=2mm,
  %frametitlerule=false,
  %repeatframetitle=false,
 % outerlinewidth=6pt,outerlinecolor=blue!20,
 pstricksappsetting={\addtopsstyle{mdfouterlinestyle}{linestyle=dashed}},
 %firstextra={\node[xshift=5cm,right,draw=White, line width=1.5pt,%star,star points=10,star point ratio=1.2, 
 %minimum size=15mm, fill=white] at (P-|O) {\tikzpicture\draw[draw=black!20,decoration=snake, fill=white] (0,0) --(2.5,0) decorate{--(2.5,1)}--(0,1)decorate{--cycle};%\includegraphics[width=10mm]{.pdf}
 %  \endtikzpicture
 %}; },
 %singleextra={\node[xshift=5cm,right,draw=White, line width=1.5pt,%star,star points=10,star point ratio=1.2, minimum size=15mm,
  %fill=white] at (P-|O) {\tikzpicture\draw[draw=black!20,decoration=snake, fill=white] (0,0) --(2.5,0) decorate{--(2.5,1)}--(0,1)decorate{--cycle};%\includegraphics[width=10mm]{.pdf}    
 %  \endtikzpicture
 %}; }
}
%%%%%%%%%%%%%%%%%%%%%%%%%%%%%%%%%%%%%%%%%%%%%%%%%%%%%%%%%%%%%%
%
%Estilo para Citar
%
%%%%%%%%%%%%%%%%%%%%%%%%%%%%%%%%%%%%%%%%%%%%%%%%%%%%%%%%%%%%%%
\mdfdefinestyle{citar}{
userdefinedwidth=.83\textwidth ,align=center,
skipabove=5pt, skipbelow=5pt,
  everyline=true,
  ignorelastdescenders=true,
  middlelinewidth=1pt,
  linecolor=blue!60,
  outerlinewidth=2pt,
  %leftline=false,
  topline=false,
  rightline=false,
  bottomline=false,
  %backgroundcolor=black!20,
  innerleftmargin=5mm, innerrightmargin=5mm,
  innerbottommargin=3mm,
  innertopmargin=1mm,
  leftmargin= 14mm,
  %frametitleaboveskip=2mm,
  %frametitlebelowskip=2mm,
  %frametitlerule=false,
  %repeatframetitle=false,
 % outerlinewidth=6pt,outerlinecolor=blue!20,
 pstricksappsetting={\addtopsstyle{mdfouterlinestyle}{linestyle=dashed}},
 %firstextra={\node[xshift=5cm,right,draw=White, line width=1.5pt,%star,star points=10,star point ratio=1.2, 
 %minimum size=15mm, fill=white] at (P-|O) {\tikzpicture\draw[draw=black!20,decoration=snake, fill=white] (0,0) --(2.5,0) decorate{--(2.5,1)}--(0,1)decorate{--cycle};%\includegraphics[width=10mm]{.pdf}
 %  \endtikzpicture
 %}; },
 %singleextra={\node[xshift=5cm,right,draw=White, line width=1.5pt,%star,star points=10,star point ratio=1.2, minimum size=15mm,
  %fill=white] at (P-|O) {\tikzpicture\draw[draw=black!20,decoration=snake, fill=white] (0,0) --(2.5,0) decorate{--(2.5,1)}--(0,1)decorate{--cycle};%\includegraphics[width=10mm]{.pdf}    
 %  \endtikzpicture
 %}; }
}
%%%%%%%%%%%%%%%%%%%%%%%%%%%%%%%%%%%%%%%%%%%%%%%%%%%%%%%%%%%%%%%%%%%%%%%%%%%%%%%%%%%%%%%%%%%%%%%%%%%%%%%%%%%%%%%%%%%%%%%%%%%%%%%%%%%%%%%%%%%%%%%%%%%%%%%%%%%%%%%%%%%%%%%%%%%%%%%%%%%%%%%%
%%%%%%%%%%%%%%%%%%%%%%%%%%%%%%%%%%%%%%%%%%%%%%%%%%%%%%%%%%%%%%%%%%%%%%%%%%%%%%%%%%%%%%%%%%%%%%%%%%%%%%%%%%%%%%%%%%%%%%%%%%%%%%%%%%%%%%%%%%%%%%%%%%%%%%%%%%%%%%%%%%%%%%%%%%%%%%%%%%%%%%%%%%%%%%%%%%%%%%%%%%%%%%%%%%%%%%%%%%%%%%%%%%%%%%
%%%%%%%%%%%%%%%%%%%%%%%%%%%%%%%%%%%%%%%%%%%%%%%%%%%%%%%%%%%%%%%%%%%%%%%%%%%%%%%%%%%%%%%%%%%%%%%%%%%%%%%%%%%%%%%%%%%%%%%%%%%%%%%%%%%%%%%%%%%%%%%%%%%%%%%%%%%%%%%%%%%%%%%%%%%%%%%%%%%%%%%%%%%%%%%%%%%%%%%%%%
\newcommand{\Dem}{\protect\tikz[overlay,yshift=-4pt]\protect\draw[line width=1.5pt,line cap=round,color=gray] (0,0)[out=14,in=180] to node [shift=(90:6.pt),color=black] {\large\bf Dem:} (7:1); \qquad \quad }
\DeclareRobustCommand{\michi}{{\mathpalette\irchi\relax}}
\newcommand{\irchi}[1]{\raisebox{\depth}{\large$#1\chi$}} % inner command, used by \rchi



%%%%%%%%%%%%%%%%%%%%%%%%%%%%%%%%%%%%%%%%%%%%%%%%%%%%%%%%%%%%%%%%%%%%%%%%%%%%%%%%%%%%%%%%%%%%%%%%%%%%%%%%%%%%%%%%%%%%%%%%%%%%%%%%%%%%%%%%%%%%%%%%%%%%%%%%%%%%%%%%%%%%%%%%%%%%%%%%%%%%%%%%%%%%%%%%%%%%%%%%%%%%%%%%%%%%%%%%%%%%%%%%%%%%%%%%%%%%%%%%%%%%%%%%%%%%%%%%%%%%%%%%%%%%%%%%%%%%%%%%%%%%%%%%%%%%%%%%%%%%%%%%%%%%%%%%%%%%%%%%%%%%%%%%%%%%%%%%%%%%%%%%%%%%%%%%%%%%%%%%%%%%%%%%%%%%%%%%%%%%%%%%%%%%%%%%%%%%%%%%%%%%%%%%%%%%%%%%%%%%%%%%%%%%%%%%%%%%%%%%%%%%%%%%%%%%%%%%%%%%%%%%%%%%%%%%%%%%%%%%%%%%%%%%%%%%%%%%%%%%%%%%%%%%%%%%
% comando aclaraciones  

\newcommand{\aclaracion}[2]{

\begin{center}
\definecolor{gris}{rgb}{0.5,0.5,0.5}
\begin{tikzpicture}[line cap=round,line join=round,scale=1]
%\draw [color=gris,dash pattern=on 1pt off 1pt, xstep=1cm,ystep=1cm] (0.1,0.1) grid (6.3,6.3);
%\draw[color=black] (0.,0.) -- (\linewidth,0.);
\node at (.23,{.9*#2}) [color=black,above right] {\bf #1};
\foreach \y in {1,2,...,#2}
\draw[shift={(0,{.9*\y})},color=gris,opacity=.5] (0.,0.) -- ({.99*\linewidth},0.);
\end{tikzpicture} 
\end{center}
}





%%%%%%%%%%%%%%%%%%%%%%%%%%%%%%%%%%%%%%%%%%%%%%%%%%%%%%%%%%%%%%%%%%%%%%%%%%%%%%%%%%%%%%%%%%%%%%%%%%%%%%%%%%%%%%%%%%%%%%%%%%%%%%%%%%%%%%%%%%%%%%%%%%%%%%%%%%%%%%%%%%%%%%%%%%%%%%%%%%%%%%%%%%%%%%%%%%%%%%%%%%%%%%%%%%%%%%%%%%%%%%%%%%%%%%%%%%%%%%%%%%%%%%%%%%%%%%%%%%%%%%%%%%%%%%%%%%%%%%%%%%%%%%%%%%%%%%%%%%%%%%%%%%%%%%%%%%%%%%%%%%%%%%%%%%%%%%%%%%%%%%%%%%%%%%%%%%%%%%%%%%%%%%%%%%%%%%%%%%%%%%%%%%%%%%%%%%%%%%%%%%%%%%%%%%%%%%%%%%%%%%%%%%%%%%%%%%%%%%%%%%%%%%%%%%%%%%%%%%%%%%%%%%%%%%%%%%%%%%%%%%%%%%%%%%%%%%%%%%%%%%%%%%%%%%%%
% comando respuesta 

\newcommand{\res}[1]{

\begin{center}
\definecolor{gris}{rgb}{0.5,0.5,0.5}
\begin{tikzpicture}[line cap=round,line join=round,scale=1]
%\draw [color=gris,dash pattern=on 1pt off 1pt, xstep=1cm,ystep=1cm] (0.1,0.1) grid (6.3,6.3);
%\draw[color=black] (0.,0.) -- (\linewidth,0.);
\node at (.23,{.9*#1}) [color=black,above right] {\bf respuesta:};
\foreach \y in {1,2,...,#1}
\draw[shift={(0,{.9*\y})},color=gris,opacity=.5] (0.,0.) -- ({.99*\linewidth},0.);
\end{tikzpicture} 
\end{center}
}

%%%%%%%%%%%%%%%%%%%%%%%%%%%%%%%%%%%%%%%%%%%%%%%%%%%%%%%%%%%%%%%%%%%%%%%%%%%%%%%%%%%%%%%%%%%%%%%%%%%%%%%%%%%%%%%%%%%%%%%%%%%%%%%%%%%%%%%%%%%%%%%%%%%%%%%%%%%%%%%%%%%%%%%%%%%%%%%%%%%%%%%%%%%%%%%%%%%%%%%%%%%%%%%%%%%%%%%%%%%%%%%%%%%%%%%%%%%%%%%%%%%%%%%%%%%%%%%%%%%%%%%%%%%%%%%%%%%%%%%%%%%%%%%%%%%%%%%%%%%%%%%%%%%%%%%%%%%%%%%%%%%%%%%%%%%%%%%%%%%%%%%%%%%%%%%%%%%%%%%%%%%%%%%%%%%%%%%%%%%%%%%%%%%%%%%%%%%%%%%%%%%%%%%%%%%%%%%%%%%%%%%%%%%%%%%%%%%%%%%%%%%%%%%%%%%%%%%%%%%%%%%%%%%%%%%%%%%%%%%%%%%%%%%%%%%%%%%%%%%%%%%%%%%%%%%%
% comando ejemplo para completar 

\newcommand{\ejparacompletar}[1]{

\begin{center}
\definecolor{gris}{rgb}{0.5,0.5,0.5}
\begin{tikzpicture}[line cap=round,line join=round,scale=1]
%\draw [color=gris,dash pattern=on 1pt off 1pt, xstep=1cm,ystep=1cm] (0.1,0.1) grid (6.3,6.3);
%\draw[color=black] (0.,0.) -- (\linewidth,0.);
\node at (.23,{.9*#1}) [color=black,above right] {\bf Ejemplo:};
\foreach \y in {1,2,...,#1}
\draw[shift={(0,{.9*\y})},color=gris,opacity=.5] (0.,0.) -- ({.99*\linewidth},0.);
\end{tikzpicture} 
\end{center}
}


%%%%%%%%%%%%%%%%%%%%%%%%%%%%%%%%%%%%%%%%%%%%%%%%%%%%%%%%%%%%%%%%%%%%%%%%%%%%%%%%%%%%%%%%%%%%%%%%%%%%%%%%%%%%%%%%%%%%%%%%%%%%%%%%%%%%%%%%%%%%%%%%%%%%%%%%%%%%%%%%%%%%%%%%%%%%%%%%%%%%%%%%%%%%%%%%%%%%%%%%%%%%%%%%%%%%%%%%%%%%%%%%%%%%%%%%%%%%%%%%%%%%%%%%%%%%%%%%%%%%%%%%%%%%%%%%%%%%%%%%%%%%%%%%%%%%%%%%%%%%%%%%%%%%%%%%%%%%%%%%%%%%%%%%%%%%%%%%%%%%%%%%%%%%%%%%%%%%%%%%%%%%%%%%%%%%%%%%%%%%%%%%%%%%%%%%%%%%%%%%%%%%%%%%%%%%%%%%%%%%%%%%%%%%%%%%%%%%%%%%%%%%%%%%%%%%%%%%%%%%%%%%%%%%%%%%%%%%%%%%%%%%%%%%%%%%%%%%%%%%%%%%%%%%%%%%
% comando conclusión 

\newcommand{\conclus}[1]{

\begin{center}
\definecolor{gris}{rgb}{0.5,0.5,0.5}
\begin{tikzpicture}[line cap=round,line join=round,scale=1]
%\draw [color=gris,dash pattern=on 1pt off 1pt, xstep=1cm,ystep=1cm] (0.1,0.1) grid (6.3,6.3);
%\draw[color=black] (0.,0.) -- (\linewidth,0.);
\node at (.23,{.9*#1}) [color=black,above right] {\bf Conclusión:};
\foreach \y in {1,2,...,#1}
\draw[shift={(0,{.9*\y})},color=gris,opacity=.5] (0.,0.) -- ({.99*\linewidth},0.);
\end{tikzpicture} 
\end{center}
}

%%%%%%%%%%%%%%%%%%%%%%%%%%%%%%%%%%%%%%%%%%%%%%%%%%%%%%%%%%%%%%%%%%%%%%%%%%%%%%%%%%%%%%%%%%%%%%%%%%%%%%%%%%%%%%%%%%%%%%%%%%%%%%%%%%%%%%%%%%%%%%%%%%%%%%%%%%%%%%%%%%%%%%%%%%%%%%%%%%%%%%%%%%%%%%%%%%%%%%%%%%%%%%%%%%%%%%%%%%%%%%%%%%%%%%%%%%%%%%%%%%%%%%%%%%%%%%%%%%%%%%%%%%%%%%%%%%%%%%%%%%%%%%%%%%%%%%%%%%%%%%%%%%%%%%%%%%%%%%%%%%%%%%%%%%%%%%%%%%%%%%%%%%%%%%%%%%%%%%%%%%%%%%%%%%%%%%%%%%%%%%%%%%%%%%%%%%%%%%%%%%%%%%%%%%%%%%%%%%%%%%%%%%%%%%%%%%%%%%%%%%%%%%%%%%%%%%%%%%%%%%%%%%%%%%%%%%%%%%%%%%%%%%%%%%%%%%%%%%%%%%%%%%%%%%%%
% boton calculadora 

\newcommand{\boton}[2]{
\tikzpicture{  \node[draw,color=#1,fill=#1,text=white,,rounded corners=3mm,text width=7mm,text height =3mm,align=center,midway]  at (0,0) { \large \bf #2 };
}
\endtikzpicture
}

\usepackage[framemethod=TikZ]{mdframed}%Paquete para estilos de los bordes de los ambientes Importante este paquete (mdframed) define un entorno  que trata automáticamente los saltos de página en el texto enmarcado.
            \mdfsetup{skipabove=2\baselineskip,skipbelow=2\baselineskip,frametitlefont=\sffamily\bfseries, needspace=4\baselineskip, splittopskip=1.5\baselineskip} 
            \mdfsetup{apptotikzsetting={\tikzset{mdfbackground/.append style={draw=none}}}} 
            
            
%%%%%%%%%%%%%%%%%%%%%%%%%%%%%%%%%%%%%%%%%%%%%%%%%%%%%%%%%%%%%%
%
%Estilo  para cajas
%
%%%%%%%%%%%%%%%%%%%%%%%%%%%%%%%%%%%%%%%%%%%%%%%%%%%%%%%%%%%%%%

\mdfdefinestyle{teoSty}{backgroundcolor=blue!10, linecolor=blue, linewidth=3pt,
skipabove=4pt,
skipbelow=3pt, 
innerleftmargin=3ex, innerrightmargin=3ex, innertopmargin=3mm, innerbottommargin=3mm, innermargin=15pt, outermargin=-15pt, needspace={0.2\textheight},roundcorner=10pt}

\mdfdefinestyle{ObsSty}{backgroundcolor=red!10, linecolor=red, linewidth=2pt,
skipabove=4pt,
skipbelow=3pt, 
innerleftmargin=3ex, innerrightmargin=3ex, innertopmargin=3mm, innerbottommargin=3mm, innermargin=5pt, outermargin=-5pt, needspace={0.2\textheight},roundcorner=10pt}

\mdfdefinestyle{EjerSty}{backgroundcolor=blue!5!green!5!, linecolor=blue!50!green!50!, linewidth=2pt,
skipabove=4pt,
skipbelow=3pt, 
innerleftmargin=3ex, innerrightmargin=3ex, innertopmargin=3mm, innerbottommargin=3mm, innermargin=5pt, outermargin=-5pt, needspace={0.2\textheight},roundcorner=10pt}



%%%%%%%%%%%%%%%%%%%%%%%%%%%%%%%%%%%%%%%%%%%%%%%%%%%%%%%%%%%%%%
%
%Estilo  para lineas de relleno usando el paquete patterns y la libreria  patterns.meta
%
%%%%%%%%%%%%%%%%%%%%%%%%%%%%%%%%%%%%%%%%%%%%%%%%%%%%%%%%%%%%%%

\tikzdeclarepattern{
  name=mylines,
  parameters={
      \pgfkeysvalueof{/pgf/pattern keys/size},
      \pgfkeysvalueof{/pgf/pattern keys/angle},
      \pgfkeysvalueof{/pgf/pattern keys/line width},
  },
  bounding box={
    (0,-0.5*\pgfkeysvalueof{/pgf/pattern keys/line width}) and
    (\pgfkeysvalueof{/pgf/pattern keys/size},
0.5*\pgfkeysvalueof{/pgf/pattern keys/line width})},
  tile size={(\pgfkeysvalueof{/pgf/pattern keys/size},
\pgfkeysvalueof{/pgf/pattern keys/size})},
  tile transformation={rotate=\pgfkeysvalueof{/pgf/pattern keys/angle}},
  defaults={
    size/.initial=5pt,
    angle/.initial=45,
    line width/.initial=.4pt,
  },
  code={
      \draw [line width=\pgfkeysvalueof{/pgf/pattern keys/line width}]
        (0,0) -- (\pgfkeysvalueof{/pgf/pattern keys/size},0);
  },
}





\newcommand{\semejante}[1]{\stackrel[#1]{}{\longleftrightarrow}}% \thicksim para decir la Semejanza de matrices la variable #1  indica la operacion elemental 
\newcommand{\OptipoUno}[1]{\stackrel[#1]{}{=}}



\newcommand{\indicador}[4]{
{\tikz[remember picture,overlay]\coordinate (#1) at #2;}{\tikz[remember picture,overlay]\coordinate (#3) at #4;} 
}

\newcommand{\marka}[2]{
{\tikz[remember picture,overlay]\coordinate (#1) at #2;}
}




\newcommand{\cof}{ \mbox{cof}}

\newcommand{\paralela}{\rotatebox[origin=c]{-10}{\Large$\parallel$}}

\usepackage{fancyhdr}
\fancypagestyle{miestilo}{
\fancyhead[L]{}
\fancyhead[C]{}
\fancyhead[R]{}
\fancyfoot[L]{}
\fancyfoot[C]{}
\fancyfoot[R]{pág. \thepage}
\renewcommand{\headrulewidth}{0pt}
\renewcommand{\footrulewidth}{0pt}
}
%%%%%%%%%%%%
%%%%%%%%%%%%
\renewcommand{\footrulewidth}{1pt}
\pagestyle{fancy}\markboth{}{}
\renewcommand{\headrulewidth}{.05pt}
\renewcommand{\headrule}{{\color{gray}%
\hrule width\headwidth height\headrulewidth \vskip-\headrulewidth}}
\renewcommand{\footrule}{{\color{gray}%
\vskip-\footruleskip\vskip-\footrulewidth
\hrule width\headwidth height\footrulewidth\vskip\footruleskip}}
\renewcommand{\footrulewidth}{.50pt}
\fancyhead[L]{\color{gray}{\footnotesize \sc  \materia} }
\fancyhead[C]{}
\fancyhead[R]{\color{gray}{\sc \footnotesize  Universidad Nacional Del Comahue}}
\fancyfoot[R]{pág. \thepage}
\fancyfoot[C]{}
\setlength{\headheight}{15pt}

%\usepackage{wrapfig}

\newcommand{\co}{\mathbb{C}}
%\newcommand{\re}{\mathbb{R}}
\newcommand{\q}{\mathbb{Q}}
\newcommand{\z}{\mathbb{Z}}
\newcommand{\n}{\mathbb{N}}
\newcommand{\ve}{\mathbb{V}}
\newcommand{\w}{\mathbb{W}}
\newcommand{\be}{{\cal B}}
\newcommand{\ca}{{\cal C}}
\newcommand{\pqx}{\mathbb{Q}[x]}
\newcommand{\prx}{\mathbb{R}[x]}
\newcommand{\pcx}{\mathbb{C}[x]}
%\newcommand{\izg}{\big\langle}
%\newcommand{\rangle}{\big\rangle}
%%%%%%%%%%%%%%%%%%%%%%%%%%%%%5

%%%%%%%%%%%%%%%%%%%%%%%%%%%%%%%%%%%%5
\newcommand{\materia}{\sc Álgebra y Geometría I}%separadas con '&'
%%%%%%%%%%%%%%%%%%%%%%%%%%%%%%%%%%%%%%%
% se define el encabezado, 
% el parametro es el texto que va despues de: trabajo practico $N^\circ$... (agregar numero y titulo)
%
%%%%%%%%%%%%%%%%%%%%%%%%%%%%%%%%%\newcommand{\encabezado}[1]{%
\fancyfoot[L]{\vspace*{-3.5mm}\color{gray}{\small T. P. 4: Vectores}}
\thispagestyle{miestilo} 
\newcommand{\encabezado}[1]{%
\fancyfoot[L]{\vspace*{-3.5mm}\color{gray}{\small T. P. #1}}
%\thispagestyle{miestilo}%
%--------------------logodpto //texto // logouni------------------
%--------------------logodpto------------------

\vspace*{-18mm}\hspace*{-11mm} \begin{minipage}{2.7cm}
\begin{center} 
     \includegraphics[width=2.7cm, height=2.7cm]{logodpto}
\end{center}  
\end{minipage}\hspace*{-15mm} \begin{minipage}{15.9cm}
\begin{center}
\begin{tabular}{c @{ - } c}
 \materia \\
\carreras
\end{tabular}

\cuatri

\medskip

TRABAJO PRÁCTICO N$^\circ$#1 

\end{center}

\end{minipage}\hspace{-14mm} \begin{minipage}{2.5cm}
  \begin{center} 
    \includegraphics[width=2.5cm, height=2.5cm]{logouni}
  \end{center}  
\end{minipage}
%-----------------------------------------------------

%----------------------------------------------------------------
\hrule
%---------------------------------------------------------------
}%%



\pagestyle{fancy}
\def\identacion{-2.5mm}%mueve la identacion de la enumeracion de los ejercicios.


\usepackage{asypictureB}
\begin{asyheader}
settings.outformat="png";
settings.render=8;
import graph3;
import smoothcontour3;

pen bpen = rgb(0.75, 0.07, 0.01);
//material m = material(diffusepen=0.7bpen, ambientpen=bpen+opacity(0.8), emissivepen=0.3*bpen,specularpen=0.999white, shininess=1.0);      
\end{asyheader}


\usepackage{multicol}
\setlength{\columnsep}{7mm}




\begin{document}


\setlength{\headheight}{15pt}

\vspace*{-2.5cm}
%--------------------logodpto------------------
\begin{figure}[htbp]
  \begin{minipage}[r][][t]{0.15\textwidth}
\begin{center} 
 % Universidad Nacional del Comahue.  
 \includegraphics[width=2.2cm, height=2.2cm]{imagen/inglogo.jpg}
\vspace{1.8cm}
\end{center}  
\end{minipage}\quad
% texto central
\begin{minipage}[r][][t]{0.31\textwidth}
\begin{center}
\vspace*{-25mm}
{\bf \materia }\\[2mm]
{\bf Módulo 4}
\end{center}
\end{minipage}\quad\begin{minipage}[r][][t]{0.31\textwidth}
\begin{center}
\vspace*{-25mm}
\textit{Ingeniería en Petróleo} \\[2mm]
Primer Cuatrimestre 2026
\end{center}
\end{minipage}\quad
%----------------------logouni-------------------------
\begin{minipage}[r][][t]{0.15\textwidth}
  \begin{center} 
     \includegraphics[width=2.2cm, height=2.2cm]{imagen/Logo_UNco.jpg}
     \vspace{1.8cm}
  \end{center}  
\end{minipage}
%-----------------------------------------------------
\end{figure}
\vspace*{-33mm}

\begin{center}
{\bf \Large{ Trabajo Práctico 4: Vectores. }}
\end{center}



%{\bf{Importante: Recordar que hay ejemplos resueltos tanto en teoría como en el archivo ejercicios modelo que te ayudarán a desarrollar este trabajo práctico.}}


\begin{enumerate}[\hspace{-3mm}$1)$]

\item  Dados el vector $\overrightarrow{PQ}$ y los puntos $A$ y $B$, en el siguiente sistema de ejes cartesianos, encontrar gráficamente:







\begin{center}
    



\definecolor{rojo}{rgb}{0.7,0.4,0.4}
\definecolor{azul}{rgb}{.4,.4,0.9}
 
\definecolor{gris}{rgb}{0.5,0.5,0.5}
\begin{tikzpicture}[line cap=round,line join=round,>=triangle 45,scale=.9]

% extremos para los ejes
\xdef\ext{6}
\xdef\f{-2}
\pgfmathtruncatemacro{\c}{\f+1}
\draw[->] (\f-.3,0) -- (12+.3,0) ;

\draw[->] (0,\f-.3) -- (0,\ext+.3) ;
\draw [color=gris!30!,dash pattern=on 1pt off 1pt, xstep=1cm,ystep=1cm] (\f-.3,\f-.3) grid (12.3,\ext+.3);

\foreach \x in {\f,...,-1,1,2,...,12}
\draw[shift={(\x,0)},color=black] (0pt,2pt) -- (0pt,-2pt) node[below] {\tiny $\x$};

\foreach \y in {\f,...,-1,1,2,...,\ext }
\draw[shift={(0,\y)},color=black] (2pt,0pt) -- (-2pt,0pt) node[left] {\tiny $\y$};


%\clip(-3.2,-1.3) rectangle (3.45,4.3);

\coordinate (A) at (6,-2);
\coordinate (B) at (2,4);
\coordinate (P) at (4,2);
\coordinate (Q) at (7,4); 
\xdef\ang{80}


%%%%%Grafico de los puntos
\draw[fill=black] (P) circle (1.5pt) node[shift=(-\ang: 5pt),right] {\scriptsize $P=(-2,1)$};

\draw[fill=black] (Q) circle (1.5pt) node[shift=(\ang: 5pt),right] {\scriptsize $Q=(1,1)$};

\draw[fill=black] (A) circle (1.5pt) node[shift=(\ang: 5pt),right] {\scriptsize $A=(-1,2)$};

\draw[fill=black] (B) circle (1.5pt) node[shift=(\ang: 5pt),right] {\scriptsize $B=(3,-1)$};


%vectores
\draw [->,line width=1.5pt, color=blue] (P) -- (Q);



\end{tikzpicture}
  
\end{center}  


   \begin{itemize}
        \item[a)] Un vector $\vec{v}$ equivalente al vector $\overrightarrow{PQ}$ con origen en el punto $(0, 0)$.
        \item[b)] Un vector, con origen en el punto $A$, que tenga igual dirección y el doble de longitud que el vector $\vec{v}$. ¿Es único? ¿Por qué?
        \item[c)] Un vector $\vec{z}$, con origen en el punto $B$, de igual dirección, sentido opuesto y la mitad de la longitud que $\vec{v}$.
        \item[d)] Un vector simétrico al vector $\vec{v}$ con respecto al eje $x$ y otro vector simétrico con respecto al eje $y$.
    \end{itemize}






\item Dados los vectores pertenecientes a $\mathbb{R}^2$, $\vec{u} = (6, 4)$, $\vec{v} = (-4, 2)$ y $\vec{w} = (-3, -2)$:

    \begin{enumerate}[a)] 
        \item Representarlos en un mismo sistema de ejes cartesianos.
        \item Calcular analítica y gráficamente las siguientes operaciones:
            \begin{enumerate}[i)]
            \begin{multicols}{3}            
                \item $\vec{w} + 2\vec{v}$.
                \item $-\tfrac{1}{2}\vec{u} + \vec{v}$.
                \item $\tfrac{1}{2}(\vec{v} + \vec{u})$.
                \item $-\vec{v} + \vec{w}$.
                \item $\tfrac{1}{2}\vec{v} + \tfrac{1}{2}\vec{u}$.
                \item $\vec{u} + 3\vec{v} - \vec{w}$.
            \end{multicols}
            \end{enumerate}
        \item ¿Qué propiedad relaciona los incisos {\scriptsize$III)$ }y {\scriptsize$V)$}?
    \end{enumerate}









\item  Dados los vectores $\vec{v} = (4, -8, 1)$, $\vec{w} = \left(\frac{1}{10}, 0, -1\right)$ y $\vec{z} = (-4, 2, 0)$:

    \begin{enumerate}[a)]
        \item Calcular analíticamente cada una de las operaciones que se indican:
            \begin{enumerate}[i)]
                \item $3(\vec{v} - 2\vec{w})$.
                \item $-2\left(\frac{1}{2}\vec{z} + \vec{v}\right) - 3\vec{w}$.
            \end{enumerate}
        \item Determinar las componentes del vector $\vec{u}$ que satisfaga: $2\vec{v} - \vec{u} + 3\vec{z} = \vec{w} - 5\vec{u}$.
    \end{enumerate}
















\item Determinar la respuesta correcta: Los siete vectores de la figura unen el centro de una circunferencia con algún punto de la misma. Si sumamos estos siete vectores el resultado:




\begin{minipage}{.6\textwidth}
\begin{center}
\begin{tikzpicture}[,>=triangle 45,rotate=-30]
    % Círculo exterior
    \xdef\r{2.1}
    \draw[dashed] (0,0) circle(\r cm);
    
    % Vectores y etiquetas
    \foreach \angle/\label in {0/$\vec{t}$, 90/$\vec{m}$, 135/$\vec{n}$, 180/$\vec{u}$, 225/$\vec{p}$, 270/$\vec{q}$, 315/$\vec{r}$} {
        \draw[->, thick] (0,0) -- (\angle:\r cm) node[pos=1.15] {\label};
    }
    
    % Etiquetas de los ángulos
    \foreach \angle in {100, 152, 202, 247 , 292} {
        \draw (\angle:.55cm) node[shift=(\angle: 5pt)]{\tiny $45^\circ$};
    }
    
    % Punto central
    \fill[blue] (0,0) circle(2pt);

    %angulos
    \draw[dashed] (0,0) ++(90:0.4) arc[start angle=90, end angle=360, radius=0.4cm];

    \draw[thick] (.3,0)--(.3,.3)--(0,.3);

\end{tikzpicture}

    
\end{center}
 \end{minipage}\quad\begin{minipage}{.38\textwidth}


    \begin{itemize}
        \item[a)] Es igual a $-\vec{p}$.
        \item[b)] Es igual a $\vec{m} + \vec{q}$.
        \item[c)] Es igual a $\vec{p}$.
        \item[d)] Es igual al vector nulo.
        \item[e)] No es ninguna de las opciones anteriores.
    \end{itemize}
\end{minipage}























\item  Determinar las componentes del vector $\vec{v}$ que tiene a los puntos $P$ como origen y $Q$ como extremo en cada uno de los siguientes casos:

    \begin{enumerate}[a) ] 
    \begin{multicols}{2}
        \item  $P = (-4, 3)$ y $Q = (1, -2)$.
        \item  $P = (2, -1, 3)$ \  y \ $Q = (4, -3, -5)$.
    \end{multicols}
    \end{enumerate}






















\item  Siendo $A = (1, 3, 1)$, $B = (1, 2, 1)$ y $C = (0, 0, 2)$ puntos pertenecientes a $\mathbb{R}^3$. Hallar:

    \begin{itemize}
        \item[a)] El punto $D$ tal que $\overrightarrow{AD}$ resulte equivalente al vector $\overrightarrow{BC}$.
        \item[b)] Un punto $M$ tal que el vector $\overrightarrow{CM}$ resulte paralelo al vector $\overrightarrow{AB}$. ¿Es único el punto encontrado? ¿Por qué?
    \end{itemize}











\item  Dados los vectores $\vec{a} = \left(1, \frac{4}{3}\right)$, $\vec{b} = (-3, 2)$ y $\vec{c} = \left(\sqrt{3}, \sqrt{3}, \sqrt{3}\right)$, calcular:

    \begin{itemize}
    \begin{multicols}{4}
        \item[a)] $\|\vec{a}\|$.
        \item[b)] $\| -\vec{a}\|$.
        \item[c)] $\|\vec{b} - \vec{a}\|$.
        \item[d)] $\| -2\vec{c}\|$.
    \end{multicols}
    \end{itemize}
 








\item  Determinar todos los valores de $k \in \mathbb{R}$ para los cuales resulta:

    \begin{itemize}
    \begin{multicols}{2}
        \item[a)] $\|\vec{v}\| = 2$, siendo $\vec{v} = (1, k, 0)$.
        \item[b)] $\|\vec{v}\| = 1$, siendo $\vec{v} = k(2, 2, 1)$.
    \end{multicols}
    \end{itemize}











\item  Calcular la distancia entre los puntos $P$ y $Q$ en cada caso:

    \begin{enumerate}[a) ]
    \begin{multicols}{2}
        \item $P = (-1, -2)$ y $Q = (1, 3)$.
        \item $P = (1, -1, 2)$ y $Q = (-3, 1, -2)$.
    \end{multicols}
    \end{enumerate}






    

\item  Sean los vectores $\vec{u} = (3, 0, 1)$ y $\vec{p} = (1, -1, 0)$, hallar:

    \begin{itemize}
        \item[a)] Un vector paralelo a $\vec{p}$ que tenga norma cinco. ¿Cuántos vectores existen con estas condiciones?
        \item[b)] El vector opuesto a $\frac{1}{2}\vec{u}$.
        \item[c)] Un versor paralelo a $(\vec{u} + \vec{p})$ y con sentido contrario.
    \end{itemize}








\item  Expresar los siguientes vectores como descomposición en las direcciones de los versores canónicos correspondientes:

    \begin{itemize}
    \begin{multicols}{2}
        \item[a)] $\vec{u} = \left(-2, \frac{4}{3}\right)$.
        \item[b)] $\vec{v} = (-1, 3, -4)$.
    \end{multicols}
    \end{itemize}





















\item  Dado el vector $\vec{u}$ del siguiente gráfico. Encontrar las componentes de dicho vector en la dirección de los vectores canónicos.




\begin{enumerate}[a) ]
\begin{minipage}{.45\textwidth}

\item \, \vspace{-9mm}
\begin{center}
     
\begin{tikzpicture}[,>=triangle 45,scale=.6]
    % Configuración inicial
    \draw[->] (-1,0) -- (5,0) node[below] {$x$}; % Eje x
    \draw[->] (0,-1) -- (0,5) node[left] {$y$}; % Eje y

    % Vector u
    \draw[->, thick, blue] (0,0) -- (45:5) node[above right] {$\vec{u}$};

    % Arco del ángulo
    \draw[thick,->] (1.5,0) arc[start angle=0, end angle=45, radius=1.5cm] node[pos=0.3,shift={({45*.5}:5mm)} ]  {\small $45^\circ$}; 

    % Magnitud del vector
    \node[left] at (-0.5,2) {\small $\|\vec{u}\|=5$};
\end{tikzpicture}


\end{center}
%------------------%---------------------%---------%------------------%---------------------%---------%------------------%---------------------%---------%------------------%---------------------%---------%------------------%---------------------%---------
\end{minipage}\qquad\begin{minipage}{.45\textwidth}

\item \, \vspace{-9mm}
\begin{center}
   
\begin{tikzpicture}[,>=triangle 45,scale=.35]
    % Configuración inicial
    \draw[->] (-6,0) -- (4,0) node[below] {$x$}; % Eje x
    \draw[->] (0,-1) -- (0,10) node[left] {$y$}; % Eje y

    % Vector u
    \draw[->, thick, blue] (0,0) -- (120:10) node[above right] {$\vec{u}$};

    % Arco del ángulo
    \draw[thick,->] (2,0) arc[start angle=0, end angle=120, radius=2cm] node[pos=0.3,shift={({120*.5}:5mm)} ]  {\small $120^\circ$}; 

    % Magnitud del vector
    \node[right] at (60:6) {\small $\|\vec{u}\|=10$};
\end{tikzpicture}


\end{center}
\end{minipage}

\end{enumerate}























\item  Un fluido fluye a una velocidad de $5\,\text{m/s}$ en dirección noroeste ($\nwarrow$). ¿Cuál es la velocidad en las direcciones este y norte?




















\item  En un sistema de tuberías de una plataforma offshore, una fuerza de $500\,\text{N}$ actúa hacia el oeste ($\leftarrow$) y otra fuerza de $300\,\text{N}$ actúa hacia el sur ($\downarrow$). ¿Cuál es la fuerza resultante en el sistema de tuberías? (describir componentes, módulo y dirección).









\item  Un buque petrolero avanza a $10\,\text{m/s}$ hacia el este ($\rightarrow$) y se encuentra con una corriente marina que va a $3\,\text{m/s}$ hacia el sur ($\downarrow$). ¿Cuál es la velocidad resultante del buque? (Indicar componentes, módulo y dirección).










\item  En cada caso, calcular el producto escalar entre los vectores $\vec{u}$ y $\vec{v}$, sabiendo que:

    \begin{itemize}
    \begin{minipage}{.35\textwidth}
        \item[a)] $\vec{u} = (5, 1)$ y $\vec{v} = (-8, 3)$.

        
    \end{minipage}\quad\begin{minipage}{.6\textwidth}
        \item[b)] $\|\vec{u}\| = 5$, $\|\vec{v}\| = 4$ y entre ellos se forma un ángulo de $\theta = \frac{\pi}{3}$.

        
    \end{minipage}
    \end{itemize}








\item   Dados los vectores $\vec{a} = (-2, \alpha)$ y $\vec{b} = (\beta, 3)$, hallar el valor de $\alpha$ y $\beta$ sabiendo que $\vec{a} \cdot \vec{b} = 2$ y $\|\vec{a}\| = \sqrt{20}$.










\item Dado un cubo donde la longitud de su arista es de $1\,\text{m}$.

    \begin{itemize}
        \item[a)] Determinar la longitud de una de sus diagonales principales.
        \item[b)] Calcular el ángulo que forma ésta con la diagonal de su base.
    \end{itemize}









\item En cada caso, hallar el ángulo que forman los vectores $\vec{r}$ y $\vec{s}$ y hallar el vector $\operatorname{Proy}_{\vec{s}} \,\vec{r}$:

    \begin{enumerate}
    \begin{multicols}{3}
        \item $\vec{r} = (1, \sqrt{3})$ y $\vec{s} = (-2, 2\sqrt{3})$.
        \item $\vec{r} = (2, 1, 1)$ y $\vec{s} = (1, -1, 2)$.
        \item $\vec{r} = (1, 2, 1)$ y $\vec{s} = (-2, 3, -4)$.
    \end{multicols}
    \end{enumerate}






    

\item En Física se considera al {\bf trabajo ($W$)} realizado sobre una partícula por una fuerza $\vec{F}$ aplicada en un desplazamiento $\vec{r}$ como $W = \vec{F} \cdot \vec{r}$.

    \begin{itemize}
        \item[a)] Calcular el trabajo realizado por la fuerza de componentes $(1, -2, 1)$ al desplazar un sólido desde un punto $A = (3, 2, -1)$ hasta $B = (2, 3, 1)$.
        \item[b)] Si una fuerza aplicada a un sólido que se desplaza una cierta distancia realizó un trabajo nulo, ¿qué puede decir de las direcciones de los vectores que representan la fuerza y el desplazamiento?
        \item[c)] Determinar el trabajo realizado por una fuerza de $5\,\text{N}$ que actúa en la dirección del vector $\vec{u} = (1, 1)$ al mover un objeto una distancia de un metro desde $P = (0, 0)$ hasta $Q = (1, 0)$.
    \end{itemize}







\item Sea el vector $\vec{u} = (3, -4)$. Hallar otro vector:

    \begin{enumerate}[a) ]
        \item Perpendicular y con el mismo módulo que el vector dado.
        \item Ortogonal al vector dado y que tenga módulo tres.
    \end{enumerate}











    

\item Dados los vectores $\vec{u} = (1, -2, 4)$ y $\vec{v} = (-2, 4, 5)$:

    \begin{itemize}
        \item[a)] Hallar un vector perpendicular a ambos.
        \item[b)] Hallar todos los vectores perpendiculares a ambos.
        \item[c)] Hallar un vector perpendicular a ambos y de módulo igual a cinco. ¿Es único este vector?
    \end{itemize}









\item Sean los vectores $\vec{v} = (a+1, -1, 4)$ y $\vec{w} = (2b+2, -2, 8)$, determinar el o los valores de $a \in \mathbb{R}$ y $b \in \mathbb{R}$ para que los vectores $\vec{v}$ y $\vec{w}$ sean paralelos. Resolver de dos maneras diferentes.












\item Sabiendo que $\vec{u} \times \vec{v} = (1, -3, \sqrt{2})$, hallar el vector $\vec{v} \times \vec{u}$.













\item Dados los vectores $\vec{a} = \iv + x\jv - \kv$ y $\vec{b} = 2\jv + \iv$, hallar todos los valores de $x \in \mathbb{R}$ para que el área del paralelogramo determinado por ellos sea $3\,\text{u}^2$.











\item Utilizando la norma del producto vectorial entre vectores, obtener el valor de la superficie total del siguiente cuerpo, donde $a > 0$:


\begin{center}
    

\begin{asy}
    import three;

size(5cm);
currentprojection = perspective(3, 2, 1);

triple A = (0, 0, 0); // Origen
triple B = (4, 0, 0); // Punto en el eje x (4a)
triple C = (0, 1, 0); // Punto en el eje y (a)
triple D = (4, 1, 0); // Parte superior del rectángulo en la base
triple E = (0, 0, 3); // Punto en la parte superior del eje z (3a)
triple F = (4, 0, 3); // Arriba a la derecha

// Ejes
draw((0, 0, 0)--(5, 0, 0), Arrow3); label("$x$", (5, 0, 0), align=(1,0,0));
draw((0, 0, 0)--(0, 2, 0), Arrow3); label("$y$", (0, 2, 0), align=(0,1,0));
draw((0, 0, 0)--(0, 0, 4), Arrow3); label("$z$", (0, 0, 4), align=(0,0,1));

// Aristas del prisma
draw(A--B, linewidth(1.5));
draw(A--C, linewidth(1.5));
draw(A--E, linewidth(1.5));
draw(C--D, linewidth(1.5));
draw(C--E, linewidth(1.5));
draw(B--D, linewidth(1.5));
draw(D--F, linewidth(1.5));
draw(E--F, linewidth(1.5));
draw(B--F, linewidth(1.5));


// dimensiones del cuerpo
draw((4,-.2,0)--(4,.2,0));
draw((0,1,-.2)--(0,1,.2));
draw((0,-.2,3)--(0,.2,3));


// Etiquetas de los puntos
label("$4a$", B+(0,-.2,0), align=(0,-1,0));
label("$a$",  C+(-.4,0,-.3), align=(0,0,-1));
label("$3a$", E+(0,.2,0), align=(-1,0,0));
\end{asy}
\end{center}
















\item  \begin{itemize}
    \item[a)] Determinar si los puntos dados son coplanares: $A = (1, 1, 6)$, $B = (2, 3, 5)$, $C = (8, 4, 6)$, $D = (2, 1, 3)$.
    \item[b)] Dados tres vectores $\vec{u} = 2\jv - \iv + 5\kv$, $\vec{v} = \alpha\jv - \kv$ y $\vec{w} = 4\kv - \jv - \iv$, hallar el valor de $\alpha \in \mathbb{R}$ de modo tal que los vectores resulten coplanares.
\end{itemize}











\item  En cada caso, dados los vectores $\vec{u}$, $\vec{v}$ y $\vec{w}$, calcular:

\begin{itemize}
    \item[a)] El volumen del paralelepípedo que ellos determinan.
    \item[b)] El área del paralelogramo que forman los vectores $\vec{v}$ y $\vec{w}$.
    \item[c)] La altura y el vector altura correspondientes al paralelepípedo cuya base está formada por los vectores $\vec{v}$ y $\vec{w}$.
        \begin{itemize}
            \item[i)] $\vec{u} = (2, 3, -2)$, $\vec{v} = (4, -1, -1)$ y $\vec{w} = (3, 1, 2)$.

            
            \item[ii)] $\vec{u} = (1, 6, -3)$, $\vec{v} = (8, 8, -4)$ y $\vec{w} = (2, -4, 2)$. 
        \end{itemize}
\end{itemize}


   
\end{enumerate}
\end{document}